\documentclass[aps,prd,twocolumn,superscriptaddress,nofootinbib,showkeys]{revtex4-2}
\usepackage{amsmath,amssymb,hyperref}
\begin{document}
\title{SERAPIS-$E$ Audited: Corrected Reheating, a Point Prediction, and How the CMB Would Measure the Compactification}
\author{Danny Eldridge}
\email{dannyeldridge54@gmail.com}
\affiliation{Independent researcher, Oklahoma City, Oklahoma, USA}
\date{17 November 2026}
\begin{abstract}
SERAPIS-$E$ is a modular plateau model in which the plateau, the vacuum at $\rho=e^{2\pi i/3}$ and a one-loop string tilt all come from the Eisenstein series $E(\tau,2)$~\cite{P3}. We report a full audit of its derivation chain, with every step recomputed under a pre-registered lock. The core arithmetic stands. Two input errors in the reheating chain are corrected: the inflaton mass at $\rho$ is $3.17\times10^{13}$~GeV, not $2.2\times10^{13}$, and the inflaton couplings were normalised with the wrong K\"ahler weight. With the MSSM $g_*=228.75$ the reheating temperature becomes $T_R=5.5\times10^8$--$1.65\times10^{10}$~GeV and $N_*=50.07$--$51.22$. The audited range is $n_s=0.97166$--$0.97596$ and $r=(4.39$--$4.90)\times10^{-4}$, partly outside the band we froze before the audit. A leptogenesis-fixed reheating temperature then gives a point prediction, $n_s=0.97439$, $r=4.82\times10^{-4}$, $\alpha_s=-9.98\times10^{-4}$, which lies inside the frozen band. A fast slow-roll code agrees with exact Mukhanov--Sasaki spectra to $|\Delta n_s|\le1.1\times10^{-6}$ over the whole range. We measure how each input moves the observables and turn the result into a forecast. If $T_R$ is known, Simons Observatory plus SPHEREx would fix the string combination $\mathrm{Im}\,S\,\mathrm{Im}\,\rho$ to $\pm14\%$. Two side results are reported: one postulate passes (a gravitino-mass window) and one is killed (a running test separating the model's two vacua). The audited model inherits the $2.4\sigma$ preference of DESI DR2 for evolving dark energy over $\Lambda$CDM. All results after the audit are post hoc and are scored only on future data.
\end{abstract}
\keywords{cosmic inflation; string cosmology; modular invariance; Eisenstein series; reheating; leptogenesis; Mukhanov--Sasaki; spectral index; tensor-to-scalar ratio; pre-registration; reproducibility; audit}
\maketitle

\section{Introduction}
SERAPIS-$E$~\cite{P3} takes the torus shape $\tau$ of the inflaton torus in a $T^6/(\mathbb{Z}_2\times\mathbb{Z}_2)$ orientifold as the inflaton, with
\begin{equation}
\frac{V}{V_0}=\frac{\hat E}{\hat E+1}\bigl(1+C\hat E\bigr),\qquad \hat E=E(\tau,2)-E(\rho,2),
\label{eq:V}
\end{equation}
and $C=15N_{\rm eff}/(64\pi^6\,\mathrm{Im}\,S\,\mathrm{Im}\,\rho)$ the one-loop K\"ahler correction of Berg, Haack and K\"ors~\cite{BHK}. Here $\mathrm{Im}\,\rho$ is the volume modulus of the four-cycle transverse to the inflaton torus, written $\mathrm{Im}\,T_1$ in Ref.~\cite{P3}. With $K=-\ln(2\,\mathrm{Im}\,\tau)$ the canonical field along $x=\tfrac12$ is $\phi=\ln y/\sqrt2$. Its predictions $n_s=0.970$--$0.975$, $r=(4.55$--$5.12)\times10^{-4}$, $\alpha_s=-(9.4$--$10.6)\times10^{-4}$ were frozen with a public hash on 15 November 2026~\cite{SERAPIS}.

A prediction is only as good as the chain behind it. Here we re-derive every link: lattice sums, the tilt, the inflaton mass, reheating, $N_*$, the spectra, and the side sectors. Each check was written as a pre-registration file and committed, with its SHA-256 hash, before any number was computed~\cite{P4}. Fixes after a lock went into a deviations file written before any rerun, and failed expectations are reported as failed. Every number below comes from a locked script whose output is in the archive~\cite{AEGIS}. The frozen band is not changed by anything here; it is still what will be scored.

\section{What the audit confirmed}
\emph{Lattice sums and the tilt.} A direct lattice sum gives $E(\rho,2)=5.7833592$ and $E(i,2)=6.0268120$. The isotropic minimum sits at $\rho$, and the cusp approach $E(\tfrac12+iy,2)/2\zeta(4)y^2=1.000065$ at $y=30$. The BHK formula reproduces $C=(1.0$--$1.3)\times10^{-6}$ for every row of the stuck-D3 tadpole $3N+N_0+N_{\rm flux}=16$. The flux-free placement $(N,N_0)=(5,1)$ gives $N_{\rm eff}=4$. A benchmark $g_s=0.232$, $\mathrm{Im}\,S=0.485$, $\mathrm{Im}\,S\,\mathrm{Im}\,\rho=781$ gives $C=1.249\times10^{-6}$. \emph{Status: computed.}

\emph{Vacuum curvature.} With $n=1$, $V''(\phi)$ at $\rho$ is $11.567\,V_0$ and $m_\phi^2/H_*^2=34.8$, as in the original derivation. \emph{Status: computed.}

\emph{Unchanged sectors.} The following are unchanged by the corrections below:
\begin{itemize}
\item the neutrino-mass input $\Sigma m_\nu=0.05878$~eV (normal ordering, $m_1\approx0$, NuFIT~6.0~\cite{NuFIT});
\item the safety caps $n_s\le0.985$, $r\le8.3\times10^{-4}$;
\item the absence of dark radiation at $\rho$ (the light sterile states thermalise, with a margin of $\times55$--$376$);
\item matter parity.
\end{itemize}

\section{What the audit corrected}
\emph{Inflaton mass.} The derivation used $m_\phi=2.2\times10^{13}$~GeV at $\rho$. Its own ratio $m^2/H^2=34.7$ then implies $H_*=3.7\times10^{12}$~GeV. The predicted $r$ with $A_s=2.1\times10^{-9}$, however, gives $H_*=5.4\times10^{12}$~GeV and $V_0^{1/4}=4.77\times10^{15}$~GeV. The corrected mass is $m_\phi=3.17\times10^{13}$~GeV. \emph{Status: computed; error found.}

\emph{Kähler weight of the couplings.} At $\rho$ the inflaton reheats by decaying to heavy neutrinos through the flavour modulus of a $\Gamma_2$ $2\oplus1$ sector, and that modulus is the inflaton itself. Two scripts mapped the potential with $y=e^{\sqrt2\phi}$ ($n=1$) but converted the couplings with $y\sqrt{2/3}$, i.e.\ $K=-3\ln(2\,\mathrm{Im}\,\tau)$, for the same field. With $n=1$ the couplings grow by $\sqrt3$ and the widths by 3. \emph{Status: computed; error found.}

\emph{Reheating and $N_*$.} A direct rerun with both fixes reproduces the rescaled $T_R$ to $\le0.15\%$. With the standard-model $g_*=106.75$ the window is $T_R=6.6\times10^8$--$2.0\times10^{10}$~GeV. The model is supersymmetric, however, so the correct count at $T_R$ is the MSSM value $g_*=228.75$. That lowers $T_R$ by $(106.75/228.75)^{1/4}=0.8265$, giving
\begin{equation}
T_R=5.5\times10^8\text{--}1.65\times10^{10}~\text{GeV},\qquad N_*=50.07\text{--}51.22 .
\end{equation}
$N_*$ is matched with the exact end of inflation ($\epsilon_H=1$), not an approximate one. The explicit $g_*$ terms in $N_*$ cancel at fixed $T_R$; we checked this to all printed digits. The frozen band assumed $N_*=48.4$--$49.9$. \emph{Status: computed; corrected.}

\emph{Running relation.} The parameter-free running relation of Ref.~\cite{P1}, with $(p,\kappa)=(2\sqrt2,1)$, had been reported as $13$--$16\%$ off. That number belonged to a predecessor potential. For SERAPIS-$E$ the relation sits $3.3$--$4.0\%$ below exact $\alpha_s$. \emph{Status: computed; correction in the model's favour.}

\emph{Second vacuum.} The variant with its vacuum at $\tau=i$ uses the computed inflaton mass $(2.28$--$2.65)\times10^{11}$~GeV. This gives $T_R=1.3$--$5.0\times10^6$~GeV and $N_*=50.40$--$50.45$, and narrows the fitted tilt to $\epsilon\approx(1.39$--$1.44)\times10^{-4}$. \emph{Status: computed.}

\emph{Two smaller corrections.}
\begin{itemize}
\item High-luminosity LHC disappearing-track searches reach winos only up to $0.9$--$1.1$~TeV~\cite{HLLHC}. They therefore cannot close the post-hoc wino window of Ref.~\cite{P3}.
\item Two ``bridge'' null tests that had been applied to SERAPIS points belong to a predecessor model. They fail on every SERAPIS point and are not tests of this model.
\end{itemize}

\section{The audited predictions}
Scanning $C=(1.0$--$1.3)\times10^{-6}$ over the corrected $T_R$ window with exact Mukhanov--Sasaki spectra gives
\begin{align}
n_s&=0.97166\text{--}0.97596,\qquad r=(4.39\text{--}4.90)\times10^{-4},\nonumber\\
\alpha_s&=-(9.11\text{--}10.14)\times10^{-4}.
\end{align}
The higher $N_*$ moves the range up by about $0.001$ in $n_s$. One corner is outside the frozen $n_s$ ceiling, and the low-$C$, high-$T_R$ corners fall below the frozen $r$ floor. We record this tension openly. Because the audit came after the freeze, the audited range is a postulate and the frozen band is what gets scored.

\emph{A point prediction.} Leptogenesis at $\rho$ works for heavy-neutrino coupling $c\gtrsim1$: 41--64\% of the scanned light-neutrino parameter space reproduces the observed baryon asymmetry at $M_H=10^{12}$~GeV. It fails our 10\% rule for $c=0.3$, reaching only 2.2\%. Taking $c=1$ and $M_H=10^{12}$~GeV fixes $T_R=7.39\times10^8$~GeV and so $N_*=50.18$. With the string benchmark $C=1.25\times10^{-6}$, the exact spectra give
\begin{equation}
n_s=0.97439,\quad r=4.82\times10^{-4},\quad \alpha_s=-9.98\times10^{-4}.
\end{equation}
This point lies inside the frozen band, $2.5\times10^{-4}$ below its $n_s$ ceiling. It is $+2.0\sigma$ from CMB-SPA ($0.9684\pm0.0030$~\cite{SPT}) and $+0.6\sigma$ from CMB-SPA+DESI ($0.9728\pm0.0027$). Varying the leptogenesis choices moves it: $c=3$ stays inside, $M_H=10^{13}$~GeV leaves in $n_s$. \emph{Status: postulate.}

\section{Toolkit cross-checks}
The spectra above come from two independent paths.
\begin{itemize}
\item \emph{Exact path:} numerical Mukhanov--Sasaki mode integration.
\item \emph{Fast path:} third-order slow-roll flow on the exact background, from the Janus calculator~\cite{Janus}.
\end{itemize}
On nine points spanning the audited range the two agree to $|\Delta n_s|\le1.1\times10^{-6}$, $\Delta r/r\approx0.011\%$ and $|\Delta\alpha_s|\le3.3\times10^{-6}$. Second-order flow alone is off by $\Delta n_s\approx-1.6\times10^{-5}$.

Getting $\alpha_s$ right from the exact path needed one fix. Starting modes at $k/aH=200$ leaves a ripple of relative size $(aH/k)^2$ in $P(k)$, which a narrow finite difference turns into $10^{-4}$-level noise in $\alpha_s$. The fix, starting at $k/aH=4000$ and fitting nine points over $\pm0.1$ in $\ln k$, is described in a companion note~\cite{P6}.

With CAMB~2.0.5~\cite{CAMB}, the primordial-tensor $D_\ell^{BB}(\ell=80)$ is $2.7$--$3.0\times10^{-5}~\mu{\rm K}^2$ for $r=4.4$--$4.9\times10^{-4}$. Every point satisfies $r<0.036$.

\section{How each input moves the sky}
We perturbed one input at a time by a fixed amount (a ``wave'') at a baseline point ($C=1.15\times10^{-6}$, $T_R=10^7$~GeV), applied it in both directions, and measured the response. All responses were linear, with up/down asymmetry $\le3\%$. Strengths are $|\Delta n_s|$ in units of the CMB-SPA error:
\begin{itemize}
\item the tilt $C$ ($\pm10\%$): $0.394\sigma$;
\item $N_*$ ($\pm1$ e-fold): $0.343\sigma$;
\item a decade in $T_R$: $0.261\sigma$.
\end{itemize}
The tilt therefore dominates, reversing the ordering we had assumed. Inside the reheating chain, the uncertainties in the inflaton mass ($\pm15\%$), in $g_*$ (SM vs MSSM) and in the K\"ahler weight ($\times\sqrt3$) move $n_s$ by only $0.016\sigma$, $0.021\sigma$ and $0.061\sigma$. The measured derivatives at the point are
\begin{align}
\frac{\partial n_s}{\partial\ln C}&=1.171\times10^{-2}, & \frac{\partial\ln r}{\partial\ln C}&=0.283,\nonumber\\
\frac{\partial n_s}{\partial\log_{10}T_R}&=7.75\times10^{-4}, & \frac{\partial\ln r}{\partial\log_{10}T_R}&=-0.0214 .
\end{align}
The implied slopes $d\ln r/dn_s$ match the edges of the audited range to 2--8\%.

\section{The CMB as a string meter}
Since $C\propto N_{\rm eff}/(\mathrm{Im}\,S\,\mathrm{Im}\,\rho)$ and $N_{\rm eff}=4$ is fixed by the tadpole, a measured $n_s$ is a measurement of the product of the dilaton and the transverse volume modulus. A Fisher forecast with the derivatives above gives:
\begin{itemize}
\item Simons Observatory with SPHEREx ($\sigma(n_s)\approx0.0016$~\cite{SOfc}): $\mathrm{Im}\,S\,\mathrm{Im}\,\rho$ to $\pm14\%$, if $T_R$ is known.
\item With LiteBIRD's optimistic $\sigma(r)\approx5\times10^{-5}$ added and $T_R$ free: $\pm19\%$, with $T_R$ to $2.76$ decades (correlation $-0.73$).
\item At LiteBIRD's requirement $\delta r=10^{-3}$, $r$ carries no information.
\end{itemize}
For the benchmark $\mathrm{Im}\,S\,\mathrm{Im}\,\rho=780$, the inflaton torus is a long thin tube at horizon exit ($y\approx14.6$) and relaxes to the hexagonal torus at the end. \emph{Status: forecast (postulate passed its pre-stated check).}

\section{Two side postulates}
\emph{Gravitino window (passes).} With $K=-\ln(2\,\mathrm{Im}\,\tau)$, positivity of the nilpotent-sector potential bounds $m_{3/2}/H$ along the trajectory. At SERAPIS-$E$'s own field values the bound is
\begin{itemize}
\item $m_{3/2}<2.0\times10^9$~GeV at horizon exit ($y=14.6$);
\item $m_{3/2}<6.7\times10^8$~GeV twelve e-folds earlier ($y=16.6$).
\end{itemize}
This is $\times960$ above the $\lesssim7\times10^5$~GeV an anomaly-mediated wino window needs. The bound was marginal for a predecessor model that exits at $y=23$--$30$; the $1/y^2$ plateau of SERAPIS-$E$ removes the problem. Non-thermal winos from gravitino decay stay at $\Omega_{\rm nth}\le4.4\times10^{-4}$. \emph{Status: postulate; consistent.}

\emph{Running as a vacuum discriminant (killed).} We asked whether the running relation's $\kappa$ separates the two vacua ($\kappa=1$ at $\rho$, $\kappa=\tfrac12$ at $i$). It does not. At $\tau=i$ the own-$\kappa$ relation is $10.8\%$ off and the crossed one $13.2\%$. Moreover the vacua differ by only $\Delta\alpha_s=1.4\times10^{-4}$, at most $0.09\sigma$ for any planned combination. \emph{Status: killed; kept on record.}

\section{Late universe}
SERAPIS fixes $\Sigma m_\nu=0.05878$~eV and has no dynamical dark energy. Fitting DESI DR2 BAO~\cite{DESI} plus a compressed CMB likelihood with this $\Sigma m_\nu$ gives $\chi^2=15.32$ for $\Lambda$CDM. A $w_0w_a$ fit with the same $\Sigma m_\nu$ improves this by $\Delta\chi^2=-7.98$ ($w_0=-0.46$, $w_a=-1.61$), a $2.4\sigma$ preference with two extra parameters. This tension is not specific to SERAPIS; it is shared by every model with a cosmological constant. A robust bound $\Sigma m_\nu<0.0588$~eV would exclude normal ordering and SERAPIS together. \emph{Status: computed.}

\section{Tests}
\begin{itemize}
\item \emph{$n_s$.} At SO+SPHEREx precision the audited range sits $\ge3.1\sigma$ from the $N_*=60$ attractor and $2.5\sigma$ from the maximal-kination attractor. This is the decisive test around 2030.
\item \emph{$r$.} $r\approx4.4$--$4.9\times10^{-4}$ is visible only if LiteBIRD reaches its optimistic sensitivity, where it would be an $\approx9\sigma$ detection.
\item \emph{$\alpha_s$.} Not testable in the 2030s.
\item \emph{Winos.} Above $1.1$~TeV, winos need gamma-ray line searches or a higher-energy collider.
\end{itemize}

\section{Limitations}
\begin{itemize}
\item The plateau factor $\hat E/(\hat E+1)$ is still a postulate~\cite{P2,P3}.
\item The $\tau=i$ tilt is fitted.
\item The wino relic window needs the temperature-dependent $g_*(T)$ and is not re-derived here.
\item All audited numbers were produced after the frozen band and after the data they are compared with. They are post hoc and will be scored only on future data.
\end{itemize}

\begin{acknowledgments}
The derivations, code and audit were developed with the assistance of Claude (Anthropic). The author directed the work and checked the results. Every check has a pre-registration file, a SHA-256 hash committed before computing, a script and an output file in the archive~\cite{AEGIS}.
\end{acknowledgments}

\begin{thebibliography}{99}
\bibitem{P3} D.~Eldridge, ``SERAPIS-$E$: Modular Plateau Inflation at the Hexagonal Torus with a One-Loop String Tilt,'' Zenodo (2026), doi:10.5281/zenodo.24290875.
\bibitem{BHK} M.~Berg, M.~Haack and B.~K\"ors, ``String loop corrections to K\"ahler potentials in orientifolds,'' JHEP \textbf{11}, 030 (2005), arXiv:hep-th/0508043.
\bibitem{SERAPIS} D.~Eldridge, ``The Shape of the Torus Sets the Sky: SERAPIS Modular Inflation, a String-Loop Tilt and a Parameter-Free Running Test,'' v3.2, Zenodo (2026), doi:10.5281/zenodo.24289118.
\bibitem{P4} D.~Eldridge, ``Lock Before You Compute: A Pre-Registration Workflow for Theoretical Physics Done with an AI Assistant,'' Zenodo (2026), doi:10.5281/zenodo.24291026.
\bibitem{AEGIS} Reproducibility archive: github.com/dannyeldridge54/aegis (locks in field\_tests/rules, outputs in field\_tests/results).
\bibitem{NuFIT} I.~Esteban \emph{et al.}, ``NuFIT-6.0,'' arXiv:2410.05380.
\bibitem{P1} D.~Eldridge, ``A Parameter-Free Consistency Relation for Plateau Inflation with an Uplift Ramp,'' Zenodo (2026), doi:10.5281/zenodo.24290583.
\bibitem{HLLHC} X.~Cid~Vidal \emph{et al.}, ``Report from Working Group 3: Beyond the Standard Model Physics at the HL-LHC and HE-LHC,'' arXiv:1812.07831.
\bibitem{SPT} SPT-3G Collaboration, arXiv:2506.20707.
\bibitem{Janus} D.~Eldridge, ``Janus String Calculator,'' v1.2, Zenodo (2026), doi:10.5281/zenodo.24299624.
\bibitem{P6} D.~Eldridge, ``A Start-Condition Ripple in Exact Mukhanov--Sasaki Running,'' companion note (2026).
\bibitem{CAMB} A.~Lewis, A.~Challinor and A.~Lasenby, Astrophys. J. \textbf{538}, 473 (2000), arXiv:astro-ph/9911177.
\bibitem{SOfc} B.~Bahr-Kalus, D.~Parkinson and R.~Easther, ``Constraining Cosmic Inflation with Observations: Prospects for 2030,'' arXiv:2212.04115.
\bibitem{DESI} DESI Collaboration, ``DESI DR2 Results II,'' arXiv:2503.14738.
\bibitem{P2} D.~Eldridge, ``Three No-Go Results for String Origins of Modular Plateau Inflation,'' Zenodo (2026), doi:10.5281/zenodo.24290738.
\end{thebibliography}
\end{document}
